\PrelimHeading{Abstract}
\PrelimEntry{Abstract}

Krishi Vaidya is an integrated crop-health diagnosis-support system developed as a final-year
B.Sc. CSIT project. The project combines a YOLO-based object-detection model, a Qwen-based
vision-language model, an Express/Mongoose backend API, and a React Native/Expo mobile
application into a single diagnosis workflow. The system is designed for image-based crop
analysis where a farmer captures or selects a crop image, submits it through the mobile app, and
receives a structured diagnosis and advisory response through backend and model-service
orchestration.

The YOLO component was trained for agricultural plant-part localization using a curated dataset
of 30,981 images and 147,467 object annotations across eleven classes. The VLM component was
prepared for disease-oriented interpretation using structured image-text data, LoRA-based
fine-tuning, quantized training, and prediction parsing. The backend implements authentication,
domain models, repositories, scan handling, diagnosis orchestration, validation middleware, and
advisory response support. The mobile app implements scan workflow, crop records, local SQLite
persistence, synchronization structure, localization, and farmer-facing result display.

The completed work demonstrates a functional integrated prototype and provides manuscript-ready
evidence through model metrics, validation outputs, generated architecture diagrams, source-backed
tables, and documented test coverage. The strongest evidence is available for the YOLO model,
VLM data/training pipeline, backend implementation, mobile architecture, and integration design.
The main limitations are the current VLM prediction reliability, missing live end-to-end runtime
evidence, incomplete mobile device validation, and the need for agricultural expert review before
real-world advisory use.

\vspace{2\baselineskip}
\noindent\textit{\fontsize{10}{12}\selectfont Keywords:} {\fontsize{10}{12}\selectfont Krishi Vaidya, crop disease diagnosis, YOLO, vision-language model, mobile agriculture, backend API}
